\begin{proof}

\begin{enumerate}[label={$\langle1\rangle$\arabic*.},leftmargin=*]
\item Choose a very general smooth fibre $F$ of $g_0$, avoiding the
images of all vertical components of $E_0$ and the jumping
sets for the pluricanonical direct images below.  Its normal
bundle in $X_0$ is trivial.  Adjunction gives
$K_F=K_{X_0}|_F$, and $c_1(K_F)=0$ integrally because
$H^2(X_0,\mathbb Z)=0$.  In particular $F$ is minimal: a
$(-1)$-curve would have canonical intersection $-1$.
Restriction of the root identity in \ref{push:src:E.2.5} gives
\[
 K_F^{A}\simeq\mathcal O_F(D E_0|_F).
\]
The right hand side is effective since $D>0$, so
$\kappa(F)\ge0$.

\item We next show $\kappa(F)\le0$, without assuming that the fibres
of an algebraic reduction have algebraic dimension zero.
If $\kappa(F)>0$, some $r>0$ has $h^0(F,K_F^r)\ge2$.
Our very-general choice ensures that this is the generic rank
of the coherent sheaf $g_{0*}K_{X_0}^r$.  A sufficiently
positive twist of this sheaf on $\mathbb P^1$ is globally
generated.  Two of its global sections can therefore be
chosen to be linearly independent on a general fibre.  By
projection formula they are sections of
$K_{X_0}^r\otimes g_0^*\mathcal O(N)$; their ratio is a
meromorphic function on $X_0$ which is nonconstant on $F$.
This contradicts the algebraic-reduction assertion of \ref{push:src:E.2.5}.
The choice of a very general fibre is legitimate: for each
$r$ the exceptional set for the generic cohomology rank is
a proper analytic subset of the base curve, and there are
only countably many $r$.

\item Thus $\kappa(F)=0$.  The classification of minimal compact
complex surfaces makes $K_F$ torsion; the relevant consequence
is also explicitly recorded in the proof of
\citet[Proposition~3.1.2]{AbashevaDeev2024AB}.
An effective divisor whose line bundle is torsion is zero:
raising its canonical section to a torsion-killing power
gives a holomorphic function on a compact connected manifold,
and that function is a nonzero constant.  Therefore
$E_0|_F=0$.  Every horizontal component of $E_0$ would meet
such a general fibre in a nonempty effective divisor.
Consequently $E_0$ has only vertical components, supported
over finitely many points of $\mathbb P^1$.

\item We now use connectedness of actual divisor limits to remove
both this vertical fixed part and the possible degree $e$.
Choose a linear combination $s=as_0+bs_1$ for which its
moving central divisor is the union of the $e$ distinct
smooth fibres over $R^{-1}([-b:a])$, all disjoint from the
support of $E_0$.  Also choose it to be a regular fibre on
one specified noncentral member.  Both requirements omit
only finitely many choices in the projective line of
combinations.  Extend it by the corresponding combination
of the frame of $\pi_*\mathcal L$.

\item Its total zero divisor $\mathcal D\subset\mathcal X$ is a
relative effective Cartier divisor: the section is not
identically zero on any fibre, by independence of the two
sections, including at zero.  In local coordinates, $t$ is
a non-zero-divisor modulo its defining equation.  Hence
$\mathcal D\to\Delta$ is proper and flat.  Its general
noncentral fibre is a smooth connected torus.  There are no
vertical irreducible components of $\mathcal D$, since such
a component would have to be the whole irreducible
threefold $X_t$.  It follows that $\mathcal D$ is reduced
and irreducible: its unique generic component has multiplicity
one, and a Cartier divisor in a smooth manifold has no
embedded components.  Normalize $\mathcal D$ and take the
Stein factorization over $\Delta$.  The finite factor has
degree one over the punctured generic locus; normality of
the disc makes it the disc itself.  Its central fibre, and
therefore the image $\mathcal D_0$, is connected.

\item But the support of
\[
 \mathcal D_0=E_0+g_0^*R^*([-b:a])
\]
is a disjoint union of the $e$ chosen smooth connected fibres
and, if nonempty, the support of $E_0$ over other base points.
Connectedness forces $E_0=0$ and $e=1$.

\item The central two sections now generate $\mathcal L_0$ at
every point.  The noncentral two sections already generate
their pencils, so the evaluation
$\pi^*\pi_*\mathcal L\to\mathcal L$ is surjective after
shrinking.  It defines the asserted map $g$.  It is proper
because $\pi$ is proper.  On each $X_t$ the nonconstant
map to the smooth curve $\mathbb P^1$ is flat: its local
structure sheaf is torsion-free over the local discrete
valuation ring of the curve.  The total source is flat
over $t$, so the local criterion for flatness applied to
the parameter $t-t_0$ gives flatness over $(z,t)$.
Equivalently, the local fibre equations form a regular
sequence in the smooth source.  At a regular
point of the central base curve, properness makes $g$ a
proper submersion on a neighbourhood of the corresponding
compact fibre.  Ehresmann's theorem identifies that central
surface differentiably with nearby complex two-tori.
The classification of compact complex surfaces, or its
special case for the differentiable four-torus, then makes
the central surface a complex two-torus as well.

\item Suppose now two positive-slope families have the same
central manifold $X_0$.  Its algebraic reduction is intrinsic
up to $\operatorname{PGL}_2$, so the preceding result makes
both central pencil lines the same non-torsion line
$L=g_0^*\mathcal O(1)$.  With indices $i=1,2$ the identities
$K_{X_0}^{A_i}=L^{D_i}$ give
$L^{D_1A_2-D_2A_1}=\mathcal O$.  The line $L$ is non-torsion
because a positive power has a nonzero section with zeros.
Thus $D_1/A_1=D_2/A_2$, and the exact recovery of the ordered
pair from this slope in \SeriesRef{push:src:E.2.3} gives the assertion.

\item Finally, a common smooth limit from opposite signs would
have a nonzero section of a positive power of $K_{X_0}$
and of a negative power, by upper semicontinuity (or by
the same injective direct-image argument).  Raising them
to a common power, their nonzero product is a constant;
both are nowhere zero.  Hence $K_{X_0}$ is torsion.
The exponential sequence and $H^1(X_0,\mathbb Z)=
H^2(X_0,\mathbb Z)=0$ make $\operatorname{Pic}(X_0)$
torsion-free, so $K_{X_0}$ would be trivial.  Its nonzero
holomorphic three-form $\Omega$ is closed and exact because
$H^3(X_0,\mathbb C)=0$.  Stokes' theorem would give
$\int_{X_0}\Omega\wedge\overline\Omega=0$, whereas its
appropriate constant multiple is a positive volume form.
This contradiction excludes the common limit.

\end{enumerate}
\end{proof}
